% =====================================================================
%  ch02b.tex —— §2.2 等价的例子：flop 双有理变换（原书 p.538–544 前）
%  底本：Orlov, Russian Math. Surveys 58:3 (2003) 511–591
% =====================================================================
% [原文 p.538]
\section{等价的例子：flop 双有理变换}

\takeaway{本节给出一整类凝聚层导出范畴等价的簇对，它们由 flop 双有理变换相联系。工具是例外族与半正交分解：先由定理 2.2.8 把爆破的导出范畴拆成 $X$ 与 $Y$ 的碎片，再证明 $\Rf\pi_*(\Lf\pi^{+*}(\,\cdot\,)\otimes\cL)$ 全忠实，并在 $l=k$ 时为等价。}

本节给出一整类光滑簇对，它们的凝聚层导出范畴彼此等价。这类例子当然早已为人所知（第一个例子是 Mukai [29] 所考虑的阿贝尔簇与其对偶簇）。与本节处理的例子相比，主要区别在于：这里得到的簇对（一般并不同构）由一个 flop 双有理变换相联系。由我们的例子还可看出，定理 2.1.3 中关于（反）典范层的条件不能减弱。

本节开头先回忆可容许子范畴与半正交分解的定义（见 [5]、[6]）。

\begin{definition}{2.2.1}
设 $\cB$ 是加性范畴 $\cA$ 的全加性子范畴。所谓 $\cB$ 在 $\cA$ 中的\emph{右正交}，指的是由一切满足如下条件的对象 $C$ 构成的全子范畴 $\cB^{\perp}\subset\cA$：对任意 $B\in\cB$ 都有 $\Hom(B,C)=0$。\emph{左正交} ${}^{\perp}\cB$ 对偶地定义。
\end{definition}

注意，若 $\cB$ 是三角范畴 $\cA$ 中的三角子范畴，则 ${}^{\perp}\cB$ 与 $\cB^{\perp}$ 也是三角子范畴。

\begin{definition}{2.2.2}
设 $I:\cN\to\cD$ 是全三角子范畴到三角范畴 $\cD$ 中的嵌入。若嵌入函子 $I$ 有右伴随 $P:\cD\to\cN$（分别地，左伴随），则称 $\cN$ 是\emph{右可容许}的（或\emph{左可容许}的）。
\end{definition}

\yizhu{可容许性的要点在于伴随的存在：右可容许意味着有右伴随 $P:\cD\to\cN$，于是每个对象都能“投影”到 $\cN$ 上，并由此得到特异三角 $N\to X\to M$，其余项落在正交补中。因此可容许子范畴像直和项一样把 $\cD$ 劈成两块；反过来，若对每个 $X$ 都能找到这样的三角，也可用它定义伴随。定理 2.2.8 中的 $D(X)$ 与诸 $D(Y)_k$ 都是这种意义下的可容许子范畴。【来源：原书 §2.2 定义 2.2.1–2.2.2 及其后一段；所引 [5]、[6]（A 级）】}

% [原文 p.539]
对子范畴 $\cN$ 而言，右可容许（或左可容许）这一性质等价于下述用正交表述的性质：对任意对象 $X\in\cD$，存在特异三角 $N\to X\to M$，其中 $N\in\cN$、$M\in\cN^{\perp}$（分别地，$M\to X\to N$，其中 $M\in{}^{\perp}\cN$、$N\in\cN$）。若一个子范畴既右可容许又左可容许，我们就简单地说它是\emph{可容许}的。

若 $\cN\subset\cD$ 是可容许子范畴，我们说 $\cD$ 具有形如 $\langle\cN^{\perp},\cN\rangle$ 或 $\langle\cN,{}^{\perp}\cN\rangle$ 的\emph{半正交分解}。这一分解过程有时可以继续下去，把子范畴 $\cN$ 或其正交进一步分解。下面给出半正交分解的一般定义。

\begin{definition}{2.2.3}
设 $(\cN_0,\dots,\cN_n)$ 是三角范畴 $\cD$ 的一列可容许子范畴。若对一切 $0\leqslant j<i\leqslant n$ 都有 $\cN_j\subset\cN_i^{\perp}$，则称这一列是\emph{半正交}的。若半正交序列生成范畴 $\cD$，即 $\cD$ 中一切包含诸 $\cN_i$ 的最小三角子范畴与 $\cD$ 重合，则称它是\emph{完备}的。此时这一序列称为范畴 $\cD$ 的\emph{半正交分解}，并记作
\[
\cD=\langle\cN_0,\dots,\cN_n\rangle.
\]
\end{definition}

\yizhu{半正交分解 $\cD=\langle\cN_0,\dots,\cN_n\rangle$ 是“带方向的直和”：条件只要求 $j<i$ 时 $\Hom(\cN_j,\cN_i)=0$，反方向的 $\Ext$ 可以非零，这正是“半”字的由来；再要求这族子范畴生成 $\cD$，每个对象便能沿方向逐层拆成特异三角。定理 2.2.8 的分解 $\langle D(Y)_{-r+1},\dots,D(Y)_{-1},D(X)\rangle$ 就是典型例子：先取 $D(X)$，再依次取扭过的 $D(Y)_k$，方向由 $k$ 的大小决定。【来源：原书 §2.2 定义 2.2.3、定理 2.2.8；所引 [5]、[6]（A 级）】}

半正交分解最简单的例子是 $\cD$ 具有完备例外族的情形。

\begin{definition}{2.2.4}
设 $E$ 是三角范畴 $\cD$ 中的对象。若 $\Hom^i(E,E)=0$（$i\neq0$）且 $\Hom(E,E)=k$，则称 $E$ 是\emph{例外}的。若一列例外对象 $(E_0,\dots,E_n)$ 生成 $\cD$，且对 $i>j$ 有 $\Hom^{\bullet}(E_i,E_j)=0$，则称它是\emph{完备例外族}。
\end{definition}

完备例外族最著名的例子由射影空间提供。

\begin{example}{2.2.5}~[2] 在射影空间 $\PP^N$ 上，对任意 $i\in\Z$，族 $(\cO(i),\dots,\cO(i+N))$ 是例外且完备的。特别地，由此得到：凝聚层的导出范畴 $\Db(\PP^N)$ 等价于有限维代数 $\End\bigl(\bigoplus_{j=0}^{N}\cO(j)\bigr)$ 上有限维模的导出范畴，而这个代数正是该例外族的自同态代数。
\end{example}

\yizhu{例外对象 $E$ 的自同态代数只有 $\Hom^0(E,E)=k$，高次 $\Ext$ 全消没，故它在导出范畴里像一点一样刚性；一族例外对象若两两之间只有单向的 $\Ext$（$i>j$ 时 $\Hom^{\bullet}(E_i,E_j)=0$）且生成整个范畴，就给出该范畴与 $\End(\bigoplus_i E_i)$ 上有限维模的导出范畴的等价。$\PP^N$ 上的 $(\cO(i),\dots,\cO(i+N))$ 是最早也最典型的例子。【来源：原书 §2.2 定义 2.2.4、例 2.2.5 及所引 [2]（A 级）】}

对其他一些簇，例如二次超曲面与旗簇，也存在类似的分解 [20]–[22]。

下面列出我们需要的关于爆破、以及凝聚层导出范畴在爆破下行为的一些事实。这些结果都包含在论文 [34] 中（另见 [7]）。设 $X$ 是光滑完备代数簇，$Y\subset X$ 是光滑嵌入的闭子簇，余维为 $r$。记 $\tilde X$ 为以 $Y$ 为中心的 $X$ 的爆破。簇 $\tilde X$ 也是光滑的，并且有交换图：
\diag{\tilde Y \arrow[r,"j"] \arrow[d,"p"'] \& \tilde X \arrow[d,"\pi"]\\
Y \arrow[r,"i"'] \& X}

% [原文 p.540]
其中 $i$ 与 $j$ 是闭嵌入，$p:\tilde Y\to Y$ 是例外除子 $\tilde Y$ 在爆破中心 $Y$ 上的射影丛；特别地，$p$ 是平坦态射。回忆 $\tilde Y\cong\PP(N_{X/Y})$，其中 $N_{X/Y}$ 是 $Y$ 在 $X$ 中的正规丛。我们用 $\cO_{\tilde Y}(1)$ 记 $\tilde Y=\PP(N_{X/Y})$ 上典范的相对丰富线丛。众所周知，该丛同构于线丛 $\cO(-\tilde Y)$ 在 $\tilde Y$ 上的限制。

% [Tutor Note: 页图此处印作 “the projective bundle of the exceptional divisor of $\tilde Y$ over the centre $Y$”，字面上是“例外除子 $\tilde Y$ 的射影丛”；按上下文，$\tilde Y$ 本身就是爆破 $\tilde X\to X$ 的例外除子，故译作“例外除子 $\tilde Y$ 在爆破中心 $Y$ 上的射影丛”。]

\begin{proposition}{2.2.6}~[34] 导出逆像函子
\[
\Lf\pi^*:\Db(X)\longrightarrow\Db(\tilde X)\qquad\text{与}\qquad p^*:\Db(Y)\longrightarrow\Db(\tilde Y)
\]
都是全忠实的。
\end{proposition}

\begin{proof}
投影公式 (3) 给出同构
\[
\Hom(\Lf\pi^*F,\Lf\pi^*G)\cong\Hom(F,\Rf\pi_*\Lf\pi^*G)\cong\Hom(F,\Rf\pi_*\cO_{\tilde Y}\otL G)
\]
（$F,G\in\Db(X)$）。对 $p^*$ 类似。把这些与 $\Rf\pi_*\cO_{\tilde X}\cong\cO_X$、$\Rf p_*\cO_{\tilde Y}=\cO_Y$ 结合起来即得证。
\end{proof}

% [Tutor Note: 页图此式末项印作 $\Rf\pi_*\cO_{\tilde Y}\otL G$，而同一证明下一行印作 $\Rf\pi_*\cO_{\tilde X}\cong\cO_X$；按投影公式应为 $\cO_{\tilde X}$，原书此处当作 $\tilde X$。本译本照录页图，读者请按 $\tilde X$ 理解。]

\begin{proposition}{2.2.7}~([34], [7]) 对 $\tilde Y$ 上任一可逆层 $\cL$，函子
\[
\Rf j_*\bigl(\cL\otimes p^*(\,\cdot\,)\bigr):\Db(Y)\longrightarrow\Db(\tilde X)
\]
是全忠实的。
\end{proposition}

\begin{proof}
为证明该函子全忠实，只需验证定理 2.1.5 的条件 1)–2)。对任意闭点 $y\in Y$，像 $\Phi(\cO_y)$ 是映射 $p$ 相应纤维的结构层（把它看作 $\tilde X$ 上的层）。不同点上的纤维互不相交，故定理 2.1.5 的正交性条件 1) 成立。

考虑某个 $p$-纤维 $F\subset\tilde Y$ 的结构层 $\cO_F$。有同构
\[
\Hom^i(j_*\cO_F,j_*\cO_F)\cong\Hom^i(\Lf j^*j_*\cO_F,\cO_F).
\]
在导出范畴 $\Db(\tilde Y)$ 中有特异三角
\[
\cO_F\otimes\cO_{\tilde Y}(1)[1]\longrightarrow\Lf j^*j_*\cO_F\longrightarrow\cO_F,
\]
其中 $\cO_{\tilde Y}(1)$ 是 $\tilde Y$ 上的相对丰富线丛，同构于 $\cO(-\tilde Y)|_{\tilde Y}$。$F$ 的纤维是射影空间，而 $\cO_{\tilde Y}(1)$ 在 $F$ 上的限制同构于 $\cO(1)$。因此
\[
\Hom^i(\cO_F\otimes\cO_{\tilde Y}(1),\cO_F)=0
\]
对一切 $i$ 成立。于是
\[
\Hom^i(j_*\cO_F,j_*\cO_F)\cong\Hom^i(\cO_F,\cO_F).
\]
所以定理 2.1.5 的条件 2) 也成立。
\end{proof}

我们用 $D(X)$ 记 $\Db(\tilde X)$ 中由 $\Db(X)$ 在 $\Lf\pi^*$ 下的像构成的全三角子范畴，用 $D(Y)_k$ 记 $\Db(\tilde X)$ 中由 $\Db(Y)$ 在 $\Rf j_*\bigl(\cO_{\tilde Y}(k)\otimes p^*(\,\cdot\,)\bigr)$ 下的像构成的全子范畴，其中 $\cO_{\tilde Y}(k)=\cO_{\tilde Y}(1)^{\otimes k}$，而 $\cO_{\tilde Y}(1)$ 是 $\tilde Y=\PP(N_{X/Y})$ 上典范的相对丰富线丛。由命题 2.2.6 与 2.2.7 可知 $D(X)\cong\Db(X)$、$D(Y)_k\cong\Db(Y)$。

% [原文 p.541]
\begin{theorem}{2.2.8}~[34] 可容许子范畴的序列
\[
\bigl\langle D(Y)_{-r+1},\dots,D(Y)_{-1},D(X)\bigr\rangle
\]
是半正交的，并给出范畴 $\Db(\tilde X)$ 的一个半正交分解。
\end{theorem}

\yizhu{定理 2.2.8 就是通常所说的 Orlov 爆破定理：$\Db(\tilde X)$ 由 $\Db(X)$ 与若干份“扭过的” $\Db(Y)$ 拼成，份数为余维 $r$ 减一。写成半正交分解后，$\Db(\tilde X)$ 的每个对象都能沿 $D(X),D(Y)_{-r+1},\dots$ 逐层拆开，于是关于爆破的问题可化归为关于 $X$ 与 $Y$ 的问题；本节后文正是用它把对象 $\overline A$ 定位到 $D(X)^{\perp}=\langle D(Y)_{-l},\dots,D(Y)_{-1}\rangle$ 之中。【来源：原书 §2.2 定理 2.2.8 及所引 [34]、[7]（A 级）】}

这一定理用原簇 $X$ 与爆破中心 $Y$ 描述了爆破 $\tilde X$ 的导出范畴。利用这一描述，我们下面研究导出范畴在最简单的 flip 与 flop 变换下的行为。考虑如下例子。

设 $Y$ 是光滑完备代数簇 $X$ 中光滑嵌入的闭子簇，满足 $Y\cong\PP^k$，且正规丛 $N_{X/Y}\cong\cO_Y(-1)^{\oplus(l+1)}$。我们假设 $l\leqslant k$。

记 $\tilde X$ 为沿中心 $Y$ 对 $X$ 的爆破。此时例外除子 $\tilde Y$ 同构于射影空间的积 $\PP^k\times\PP^l$。此外，在这种情形下，$\tilde Y$ 在 $\tilde X$ 中的正规层有如下描述：
\[
N_{\tilde X/\tilde Y}=\cO_{\tilde X}(\tilde Y)|_{\tilde Y}\cong\cO(-1;-1),
\]
其中 $\cO(-1;-1):=p_1^*\cO_{\PP^k}(-1)\otimes p_2^*\cO_{\PP^l}(-1)$。这些事实使我们断言：存在 $\tilde X$ 的一个收缩，使 $\tilde Y$ 投影到第二个因子 $\PP^l$。该收缩存在于解析范畴中，其结果是一个光滑簇 $X^+$，它一般可能不是代数的。我们假设 $X^+$ 是代数的。上面描述的几何都包含在下图中：
\begin{equation}
\parbox{0.92\linewidth}{\centering\diag{
\& \tilde Y \arrow[dl,"p"'] \arrow[d,"j"] \arrow[dr,"p^+"] \& \\
Y \arrow[d,"i"'] \& \tilde X \arrow[dl,"\pi"'] \arrow[dr,"\pi^+"] \& Y^+ \arrow[d,"i^+"] \\
X \& \& X^+
}}
\tag{20}
\end{equation}

双有理映射 $\mathrm{fl}:X\dashrightarrow X^+$ 是 flip 或 flop 的最简单例子。当 $l<k$ 时它是 flip，当 $l=k$ 时它是 flop。

% [Tutor Note: p.541 图 (20) 的页图由缩微胶片扫描而来，底排两个节点（$X$ 与 $X^+$）的标签、以及右下角若干残段均已糊死，无法逐字辨认；页图右端残留的短虚线亦不能判定归属。本图按正文所述映射重绘：6 个节点为 $\tilde Y$、$Y$、$\tilde X$、$Y^+$、$X$、$X^+$，7 条实线箭头为 $p:\tilde Y\to Y$、$j:\tilde Y\to\tilde X$、$p^+:\tilde Y\to Y^+$、$i:Y\to X$、$\pi:\tilde X\to X$、$\pi^+:\tilde X\to X^+$、$i^+:Y^+\to X^+$。凡页图不能证实的虚线箭头（如以 $\mathrm{fl}$ 标注的 $X\dashrightarrow X^+$）一概不画，以免杜撰。]

\yizhu{flip 与 flop（翻模）都是双有理变换，都把 $X$ 中的子簇 $Y$ 换成 $X^+$ 中的另一子簇 $Y^+$，且被改动的子簇余维 $\geqslant2$。区别在典范类：flip 两侧的典范类不同，收缩曲线上的 $K$ 由负变正；flop 两侧的典范类相同，$K_{X}$ 与 $K_{X^+}$ 在公共的爆破 $\tilde X$ 上重合。本例正体现这一点：$l<k$ 时 $Y=\PP^k$ 被换成维数更低的 $Y^+=\PP^l$，是 flip；$l=k$ 时 $Y$ 与 $Y^+$ 都是 $\PP^k$，而式 (21) 中 $\omega_{\tilde X}|_{\tilde Y}=\cO(-k;-k)$ 关于 $\tilde Y=\PP^k\times\PP^k$ 的两个因子对称，是 flop。【来源：原书 §2.2（$l<k$ 与 $l=k$ 的区分）与式 (21)；flip、flop 的一般定义见 Kollár–Mori《Birational Geometry of Algebraic Varieties》（A 级）】}

以下我们需要典范层 $\omega_{\tilde X}$ 在除子 $\tilde Y$ 上的限制的公式。对光滑子簇的爆破，我们得到
\[
\omega_{\tilde X}\cong\pi^*\omega_X\otimes\cO_{\tilde X}(l\tilde Y).
\]
伴随公式给出
\[
\omega_X|_Y\cong\omega_Y\otimes\bigwedge^{l+1}N^*_{X/Y}\cong\cO_Y(l-k).
\]
把这些事实合起来，得到同构
\begin{equation}
\begin{aligned}
\omega_{\tilde X}|_{\tilde Y}&\cong\bigl(\pi^*\omega_X\otimes\cO_{\tilde X}(l\tilde Y)\bigr)\big|_{\tilde Y}\\
&\cong p^*(\omega_X|_Y)\otimes\cO_{\tilde X}(l\tilde Y)|_{\tilde Y}\cong\cO(-k;-l).
\end{aligned}
\tag{21}
\end{equation}
本节的主定理把 $X$ 与 $X^+$ 上凝聚层的导出范畴联系起来。

% [原文 p.542]
\begin{theorem}{2.2.9}
设 $\cL$ 是 $\tilde X$ 上的线丛。在上述记号下，函子
\[
\Rf\pi_*\bigl(\Lf\pi^{+*}(\,\cdot\,)\otimes\cL\bigr):\Db(X^+)\longrightarrow\Db(X)
\]
是全忠实的。
\end{theorem}

\begin{proof}
先考虑 $\cL$ 在 $\tilde Y$ 上的限制。由于 $\tilde Y=\PP^k\times\PP^l$，存在整数 $a$ 与 $b$ 使 $\cL|_{\tilde Y}\cong\cO(a;b)$。

我们要证明：对任意一对 $A,B\in\Db(X^+)$，复合映射
\begin{align}
\Hom(A,B)&\xrightarrow{\ \sim\ }\Hom(\Lf\pi^{+*}A,\Lf\pi^{+*}B)\notag\\
&\longrightarrow\Hom\bigl(\Rf\pi_*(\Lf\pi^{+*}A\otimes\cL),\Rf\pi_*(\Lf\pi^{+*}B\otimes\cL)\bigr)
\tag{22}
\end{align}
是同构。利用函子的伴随性，我们得到同构
\[
\Hom\bigl(\Rf\pi_*(\Lf\pi^{+*}A\otimes\cL),\Rf\pi_*(\Lf\pi^{+*}B\otimes\cL)\bigr)
\cong\Hom\bigl(\Lf\pi^*\Rf\pi_*(\Lf\pi^{+*}A\otimes\cL),\Lf\pi^{+*}B\otimes\cL\bigr).
\]

考虑特异三角
\begin{equation}
\Lf\pi^*\Rf\pi_*(\Lf\pi^{+*}A\otimes\cL)\longrightarrow\Lf\pi^{+*}A\otimes\cL\longrightarrow\overline A.
\tag{23}
\end{equation}

于是，要证明复合映射 (22) 是同构，必要且只需证明
\begin{equation}
\Hom(\overline A,\Lf\pi^{+*}B\otimes\cL)=0.
\tag{24}
\end{equation}

由命题 2.2.6，复合 $\Rf\pi_*\Lf\pi^*$ 同构于恒等函子；因此把函子 $\Rf\pi_*$ 作用到特异三角 (23) 上，即知 $\Rf\pi_*\overline A=0$。于是对任意对象 $C\in\Db(X^+)$ 都有 $\Hom(\Lf\pi^*C,\overline A)=0$。因此 $\overline A$ 属于子范畴 $D(X)^{\perp}$。

定理 2.2.8 给出半正交分解
\[
D(X)^{\perp}=\bigl\langle D(Y)_{-l},\dots,D(Y)_{-1}\bigr\rangle.
\]

由于 $Y$ 是射影空间，由例 2.2.5 可知每个 $D(Y)_{-i}$ 都具有完备例外族。把这些族汇集起来，就得到 $D(X)^{\perp}$ 中的一个完备例外族。下面的族对我们较为方便：
\[
D(X)^{\perp}=\left\langle
\begin{array}{l@{\quad}c@{\quad}l}
\Rf j_*\cO(a-k;-l), & \dots\quad\dots & \Rf j_*\cO(a;-l),\\
\Rf j_*\cO(a-k+1;-l+1), & \dots\quad\dots & \Rf j_*\cO(a+1;-l+1),\\
\vdots & & \vdots\\
\Rf j_*\cO(a-k+l-1;-1), & \dots\quad\dots & \Rf j_*\cO(a+l-1;-1)
\end{array}\right\rangle.
\]

% [原文 p.543]
% [Tutor Note: 文字层把 (23)、(24) 及以下各处的 $\overline A$ 一律抽成 $A$（上横线丢失）；页图三处均印作 $\overline A$。本译本从页图，全篇作 $\overline A$，以免与证明开头的对象 $A\in\Db(X^+)$ 混淆。]

现在我们可以重排这个例外序列，得到 $D(X)^{\perp}$ 的如下形式的半正交分解
\[
D(X)^{\perp}=\langle\cB,\cA\rangle,
\]
其中 $\cA$ 与 $\cB$ 分别是由 $i\geqslant a$ 与 $i<a$ 的 $\Rf j_*\cO(i;s)$ 生成的子范畴。对 $1\leqslant i\leqslant k$ 与 $1\leqslant s\leqslant l$，对象 $\Rf j_*\cO(a-i;-s)$ 同时属于子范畴 $D(X)^{\perp}$ 与 $D(X^+)^{\perp}\otimes\cL$。特别地，$\cB\subset D(X)^{\perp}\cap(D(X^+)^{\perp}\otimes\cL)$。把 $\Hom$ 作用到特异三角 (23) 上，得到
\[
\Hom(\overline A,\Rf j_*\cO(a-i;-s))=0\qquad\text{对 }1\leqslant i\leqslant k\text{ 与 }1\leqslant s\leqslant l.
\]

由于 $\overline A\in D(X)^{\perp}$ 且 $\overline A$ 与子范畴 $\cB$ 正交，立刻推得 $\overline A\in\cA$。又注意：若对象 $\Rf j_*\cO(a+i;s)$ 属于子范畴 $\cA$，则 $i$ 满足不等式 $0\leqslant i<l$。考虑到典范类 $\omega_{\tilde X}|_{\tilde Y}\cong\cO(-k;-l)$ 的公式 (21) 以及条件 $l\leqslant k$，可见 $\cA\otimes\omega_{\tilde X}\subset D(X^+)^{\perp}\otimes\cL$。因此，对任意对象 $B\in\Db(X^+)$ 都有
\[
\Hom(\Lf\pi^{+*}B\otimes\cL,\overline A\otimes\omega_{\tilde X})=0.
\]
应用 Serre 对偶 (6)，立即得到所要的等式 $\Hom(\overline A,\Lf\pi^{+*}B\otimes\cL)=0$。
\end{proof}

\begin{theorem}{2.2.10}
在上述记号下，若 $l=k$（因而 $\mathrm{fl}$ 是 flop），则函子 $\Rf\pi_*\bigl(\Lf\pi^{+*}(\,\cdot\,)\otimes\cL\bigr)$ 是三角范畴的等价。
\end{theorem}

\begin{proof}
由前一定理，所论函子全忠实。它的左伴随形如 $\Rf\pi^{+}_*\bigl(\Lf\pi^*(\,\cdot\,)\otimes\cL'\bigr)$，其中 $\cL'=\cL^{-1}\otimes\omega_{\tilde X}\otimes\pi^{+*}\omega_{X^+}^{-1}$。因此由前一定理它也是全忠实的。这就证明了两个函子都是等价。
\end{proof}

我们指出，若在某个由彼此不交、且各自满足定理条件的子簇组成的有限集 $Y_1,\dots,Y_s$ 上同时作 flop，上述断言的证明依然成立。这一简单的说明与我们关于所得簇 $X^+$ 是代数的假设密切相关。原因在于：有许多例子表明，只在其中一个子簇 $Y_i$ 上作上述类型的双有理变换，所得簇不是代数的；而同时在整个集合上作 flip（或 flop），所得簇则是代数的。

第二点评注是：定理 2.2.10 中的 flop 簇 $X$ 与 $X^+$ 当然一般并不同构；flop 在双有理几何中经常出现，例如在所谓“自一条直线作双投影”以描述 Fano 三维簇的构造中（见 [19]，§ 8）。设 $V$ 是指标为 1、$\Pic V=\Z$ 的 Fano 三维簇，它由其反典范系统嵌入射影空间。于是，把该簇沿一条直线爆破，得到簇 $X$，其反典范类“几乎”是丰富的；也就是说，由反典范系统定义的映射把该簇上某一族曲线收缩，这些曲线正是 $V$ 上与所爆破直线相交的直线在 $X$ 中的真变换。在许多例子中
% [原文 p.544]
这些曲线的正规层为 $\cO(-1)\oplus\cO(-1)$，于是我们正处在定理的情形。在这些曲线上同时作 flop，得到与 $X$ 不同构的簇 $X^+$；然而由定理 2.2.10，它具有相同的凝聚层导出范畴。特别地，这个例子表明：定理 2.1.3 中关于反典范类的丰富性条件，对于由 $\Db(X)$ 还原 $X$ 来说不能减弱。对于极小模型纲领中出现的一般型簇，也有类似的例子。

这些结果还有另一个自然的推广。设光滑完备代数簇 $X$ 中的光滑子簇 $Y$ 是光滑簇 $Z$ 上秩为 $k+1$ 的向量丛 $E$ 的射影化，即 $Y\cong\PP(E)\to Z$。我们还假设正规丛 $N_{X/Y}$ 限制到映射 $Y\to Z$ 的纤维上同构于 $\cO_{\PP^k}(-1)^{\oplus(l+1)}$。我们同样假设 $l\leqslant k$。仍记 $\tilde X$ 为以 $Y$ 为中心的 $X$ 的爆破，我们再次得到形如 (20) 的图，其中 $Y^+$ 是 $Z$ 上秩为 $l+1$ 的丛的射影化。在这种情形下可以断言：定理 2.2.9 与 2.2.10 的类似结论依然成立。

当 $X$ 是三维簇、$Y$ 是满足 $Y\cdot K_X=0$ 的有理曲线时，还会出现其他类似的例子。此时 $Y$ 上的正规丛可能形如 $\cO(-1)\oplus\cO(-1)$、$\cO\oplus\cO(-2)$ 或 $\cO(1)\oplus\cO(-3)$。在这些情形下都存在 flop 双有理变换 $\mathrm{fl}:X\dashrightarrow X^+$。此外，在每种情形下 $X$ 与 $X^+$ 的凝聚层导出范畴都等价。第一种情形是定理 2.2.10 的特例。第二种情形在 [7] 中处理。更近一些，[10] 一并证明了所有这些情形下的范畴等价。

% [Tutor Note: 本块止于此处：其下即为第 3 章标题 “CHAPTER 3 — Fully faithful functors between derived categories”，属 ch03a 的内容。]

\chapterreview{第 2 章处在全书的枢纽位置。第 1 章备好了三角范畴与导出范畴的语言，第 3 章要证明“导出范畴之间的每个全忠实函子都由积上的对象表示”；第 2 章做的正是把这两端接起来：一方面把 $\Db(X)$ 当作具体对象来研究（2.1），另一方面引入由积上对象定义的函子 $\Phi_{\cE}$ 并证明第一批真正的等价（2.2）。没有 2.1 的分次代数 $\mathrm{HA}(X)$，关于“(反)典范类丰富时范畴还原簇”的定理 2.1.3 就无从表述；没有 2.2 的 flop 例子，“何时发生等价”这个问题的分量也无从显现。

两条路线方向相反。2.1 是“不变量”路线：从 Serre 对偶与投影公式出发，把 $\bigoplus_{i,k}\mathrm{HA}^{i,k}(X)$ 与 Hodge 菱形竖线上的上同调维数做成导出等价的不变量（定理 2.1.8、推论 2.1.10），再用定理 2.1.5 的全忠实判据把“函子是否全忠实”化归为对点与纤维的正交性计算。2.2 是“构造”路线：先以例外族与半正交分解（定义 2.2.1–2.2.4）给出拆解范畴的语法，再用爆破定理 2.2.8 把 $\Db(\tilde X)$ 写成 $\Db(X)$ 与 $r-1$ 份扭过的 $\Db(Y)$ 的拼装，最后沿 flop 的两侧把这两块对上。

硬结论有三条。其一，导出等价保持 $\mathrm{HA}(X)$ 与 Hodge 菱形竖线上的维数（定理 2.1.8、推论 2.1.10），这是“范畴记得几何”的定量版本。其二，爆破的导出范畴有半正交分解 $\Db(\tilde X)=\langle D(Y)_{-r+1},\dots,D(Y)_{-1},D(X)\rangle$（定理 2.2.8），其中 $D(X)\cong\Db(X)$、$D(Y)_k\cong\Db(Y)$。其三，对 $Y\cong\PP^k$、$N_{X/Y}\cong\cO_Y(-1)^{\oplus(l+1)}$、$l\leqslant k$ 的 flop，任一 $\tilde X$ 上的线丛 $\cL$ 给出的 $\Rf\pi_*(\Lf\pi^{+*}(\,\cdot\,)\otimes\cL)$ 全忠实，$l=k$ 时还是等价（定理 2.2.9、2.2.10）。三者合起来说明：flop 两侧的导出范畴等价，而 flip 情形本节只证到全忠实、没有给出等价；(反)典范类的丰富性条件确实不能减弱。

读这一章要当心四处。第一，半正交只要求单向消没（$j<i$ 时 $\Hom(\cN_j,\cN_i)=0$），反方向允许非零，所以它不是正交直和，$\langle\cA,\cB\rangle$ 与 $\langle\cB,\cA\rangle$ 一般不同。第二，$D(Y)_k$ 不是 $\Db(Y)$ 本身，而是经 $\Rf j_*(\cO_{\tilde Y}(k)\otimes p^*(\cdot))$ 扭过的拷贝；不同 $k$ 的拷贝互相正交，这正是分解里出现 $r-1$ 份的原因。第三，定理 2.2.9 对任意线丛成立，但证明的关键一步用到式 (21) 的 $\omega_{\tilde X}|_{\tilde Y}\cong\cO(-k;-l)$ 与条件 $l\leqslant k$，缺一不可。第四，“$X^+$ 是代数的”是额外假设：收缩本身只在解析范畴中存在，正文特意补上“同时在多个不交子簇上作 flop”的说明来填这个缺口。

从今天的眼光看，2.2 的意义不止于“又一批等价的例子”。它示范了一种方法：先用半正交分解把几何操作翻译成范畴的拼装，再在两个拼装之间构造函子并验证全忠实。后来的 Lefschetz 分解（Kuznetsov）与关于 flop 不变量的许多工作，都沿这条路线展开。反过来，flop 的例子中 $X$ 与 $X^+$ 一般不同构却导出等价，也提醒读者：$\Db(X)$ 是比 $X$ 更粗的不变量，问“它能记得多少”时必须把 (反)典范类的正性放在前提里。}
